\documentclass[runningheads]{llncs}
\usepackage{graphicx}
\usepackage{booktabs}
\usepackage{xcolor}
\usepackage[utf8]{inputenc}
\usepackage{hyperref}
\renewcommand\UrlFont{\color{blue}\rmfamily}

\begin{document}

\title{The Price of a Proof: A Controlled Token-Cost Comparison of Formal
Contract-Driven Development and Ordinary Development-plus-Review}
\titlerunning{The Price of a Proof}

\author{Withheld for review}
\authorrunning{Withheld for review}
\institute{Study repository and full data available; see Data Availability}

\maketitle

\begin{abstract}
Formal methods are defended on correctness and attacked on cost, yet the cost
claim is seldom measured against a realistic alternative doing the same work.
We ran a controlled comparison on a 7{,}343-line Z80 assembly chess engine into
which an independent agent seeded seven single-byte faults, each with a sealed
acceptance test written before any repair was attempted. Two arms fixed every
fault: ordinary development-plus-review, and Formal Contract-Driven Development
(FCDD) working from a pre-built machine-checked contract. Across seven matched
pairs FCDD was more expensive on 7 of 7 (exact two-sided sign test $p=0.016$),
median ratio $3.40\times$; its upfront contract never amortised. Replicating the
pair whose arms were closest revealed that the ordinary arm is \emph{bimodal} —
it produced either a minimal one-byte repair (\$8.66, \$11.78) or a redesign
(\$44.27, \$41.35) — while FCDD produced the minimal repair in all three runs.
Matched on solution, FCDD costs a median $5.0\times$ more. Blinded, order-
counterbalanced grading of 14 judgements favoured FCDD on 6 of 7 pairs
($p=0.125$), concentrated in test quality. We report three measurement findings
we regard as more transferable than the ratio: an instruction to review ``until
no findings remain'' has no fixed point; measured cost covaries with review-loop
length; and processes killed by infrastructure emit no cost record.

\keywords{Formal methods \and Empirical software engineering \and Cost
measurement \and LLM agents \and Code review}
\end{abstract}

\section{Introduction}

The case for formal methods is argued on one axis and rebutted on another.
Proponents show what proof buys; opponents answer that it costs too much. The
disagreement persists partly because the cost side is rarely measured against a
competent alternative performing the identical repair on the identical defect,
judged by the identical gate.

This study measures it, and its more useful contribution turned out to be how
badly that measurement can go wrong. Five times a confident, plausible finding
proved to be an artefact of our own apparatus rather than a property of the
methods; three were caught by our own mechanical checks and two only by
adversarial review of a finished draft. Each correction moved the headline, and
two reversed its direction.

\section{Related Work}

\paragraph{Empirical cost of formal methods.}
The field's own guidance notes that empirical studies of formal methods are
rare~\cite{terbeek2022}. What exists is largely retrospective single-project
accounting with no control arm: seL4 required 2.2 person-years to build and over
20 to verify~\cite{staples2014,klein2009}; ShardStore's authors compare their
overhead to figures \emph{drawn from other papers}~\cite{bornholt2021}; the AWS
TLA+ reports rest on counterfactual testimony~\cite{newcombe2015}. The largest
survey concedes the gap outright: in a majority of projects surveyed there is no
data available on relative time and costs~\cite{woodcock2009}. The classic
negative industrial result is observational and unpriced~\cite{pfleeger1997}.

\paragraph{The one controlled precedent.}
Sobel and Clarkson ran a controlled formal-versus-informal experiment with
twenty student teams~\cite{sobel2002}; we therefore do not claim the first
controlled comparison. They measured product quality and never priced either
arm, and Berry and Tichy argued in print that the design could not support even
the quality claim~\cite{berry2003}. Our sealed acceptance tests and
order-counterbalanced blinded grading respond directly to that critique. The
closest work in unit of analysis reports 87 minutes to develop each of 73 unit
proofs~\cite{amusuo2025}, but against a detection-rate baseline rather than a
competing method, and in human minutes rather than tokens.

\paragraph{Agentic repair and its cost.}
LLM agents repairing real defects is established~\cite{jimenez2024}, as is the
finding that simpler structured pipelines can beat autonomous agents on quality
and dollars alike~\cite{xia2024}. Our contribution is therefore not that
structure helps but \emph{which} structure, at what price. Cost-aware agent
evaluation is likewise prior art~\cite{kapoor2024}, and per-task token
consumption has been characterised in detail, including the finding that runs of
the same task can differ by up to $30\times$~\cite{bai2026} --- a noise floor any
cost delta must clear. Benchmark contamination motivated seeding our own
faults~\cite{aleithan2024}. The nearest neighbour varies one prompt line
controlling specification grounding for $+38$ percentage points of correct
code~\cite{haeri2026}, but reports no cost accounting.

\paragraph{Review effectiveness --- our instrument findings are replication.}
That review surfaces far more maintainability issues than defects is long
established --- 75\% evolvability to 25\% functional~\cite{mantyla2009},
replicated in open source with 7--35\% of comments discarded~\cite{beller2014}.
Porter et al.\ already quantified the signal-to-noise ratio we
observed~\cite{porter1998}, and Rigby and Bird state from industrial data that
more reviewers yield ``only a minimal increase in the number of comments'' and
``no increase in the number of change sets submitted''~\cite{rigby2013}. Our 85
findings without a defect in the fix is therefore a replication. What appears
unclaimed is the \emph{stopping rule}: over 80\% of human changes involve at most
one review iteration~\cite{sadowski2018}, so twelve rounds under an unbounded
instruction is an order-of-magnitude departure.

\paragraph{Seeded faults and grading: threats we inherit.}
Mutants couple to 73\% of real faults but not to 17\%~\cite{just2014};
correlation with real-fault detection weakens once suite size is
controlled~\cite{papadakis2018}; hand-seeded faults tend to be harder to detect
than real ones~\cite{andrews2005}. Our defence is that single-statement bugs are
an abundant measured class~\cite{karampatsis2020}. Position bias in LLM judges
was reported by Zheng et al.~\cite{zheng2023}, demonstrated as rank-hacking by
reordering~\cite{wang2024}, and shown strongest when the quality gap is
small~\cite{shi2025} --- exactly our regime. Our first grading round reproduced
this failure and we applied the established remedy; that is a confession, not a
finding. For method comparison we follow standard guidance~\cite{ko2015}, and
adopt the pair-programming meta-analysis framing of reporting effort as a
separate dependent variable rather than equalising it~\cite{hannay2009}.

\paragraph{What remains ours.}
Searching for prior work pricing a formal-methods workflow in LLM token cost
against a conventional one returned nothing we could verify; we state that as a
searched-and-not-found result. The defensible contribution is the conjunction: a
concurrently-run control arm on the same faults, a directly metered cost unit, a
repeated-measures design, pre-committed outcome measurement, and the
stopping-rule and bimodality results.

\section{Method}

An independent agent seeded seven single-fault variants into the engine, each
differing from pristine by exactly one byte in the assembled binary, wrote
symptom-level reports with no location hints, and sealed an answer key whose
hash was published before any repair began. Each arm$\times$defect cell ran as an
isolated headless session on the same model at the same effort. FCDD's contract
was built once as a separately measured phase (1{,}271 lines of specification, 95
machine-checked theorems, empty axiom profile) at \$21.75.

Cost is reported three ways --- the platform figure, the same tokens priced with
our pre-registered rate table, and raw token totals --- because they disagree.
Token totals come from each session's \texttt{modelUsage}, which includes
reviewer subagents; an earlier ledger recorded main-loop usage only and
under-counted by 29--99\%.

\paragraph{Disclosed asymmetry.} FCDD's contract was built by reading the
\emph{pristine} engine, so Arm~B held an executable model of correct behaviour
covering all seven faults by construction; Arm~A received only the symptom
report. This is the most consequential design flaw and would not transfer to
faults of omission.

\section{Instrument findings}\label{sec:instr}

\textbf{Unbounded review has no fixed point.} In one cell Arm~A fixed the seeded
fault correctly on its first pass, then ran 12 review rounds producing 85
findings, none of which was a defect in the fix (OLS slope $+0.025$
findings/round, $R^2=0.001$). Bounding the rule --- a round finding no defect
closes the gate, cap three rounds --- cut the same cell from at least \$128.62
with no completion to \$7.32 with a passing gate.

\textbf{Cost covaries with review effort.} Over all 13 final-instrument cells,
cost $=4.13+21.35\cdot\mathrm{ArmB}+7.13\cdot\mathrm{reviewers}$, $R^2=0.826$;
controlling for reviewer count attenuates the arm gap from \$30.34 to \$21.35.

\textbf{Killed processes emit no cost record.} One cell reported \$1.06; its
first attempt had been killed without emitting a result record. Priced from
recovered transcripts at our pre-registered rates the true figure is \$23.98 ---
an under-count of roughly $20\times$.

\section{Results}

\begin{table}[t]
\caption{Cost per defect, final instrument. $\dagger$ imputed from recovered tokens.}
\label{tab:cost}
\centering
\begin{tabular}{lrrr}
\toprule
Defect & Arm~A & Arm~B & Ratio\\
\midrule
bug01 quiescence & \$44.27 & \$49.63 & $1.12\times$\\
bug02 fifty-move & \$7.32 & \$50.37 & $6.88\times$\\
bug03 promotion & \$11.96 & \$23.98$\dagger$ & $2.01\times$\\
bug04 eval sign & \$11.30 & \$55.40 & $4.90\times$\\
bug05 castling & \$8.88 & \$31.05 & $3.50\times$\\
bug06 TT bound & \$16.44 & \$48.69 & $2.96\times$\\
bug07 mate detect & \$7.11 & \$24.19 & $3.40\times$\\
\midrule
\multicolumn{3}{l}{Median ratio (platform / rates / tokens)} & $\mathbf{3.40\times}$ / $3.28\times$ / $4.74\times$\\
\bottomrule
\end{tabular}
\end{table}

FCDD was dearer on 7 of 7 pairs (sign test $p=0.016$), the only statistically
supportable cost result here (Table~\ref{tab:cost}, Fig.~\ref{fig:cost}). No
crossover exists: $N$ is undefined when the saving is negative on every pair.

\begin{figure}[t]
\centering
\includegraphics[width=\textwidth]{fig2_cost_per_defect.pdf}
\caption{Cost per defect. FCDD is more expensive on every pair.}
\label{fig:cost}
\end{figure}

\paragraph{Replication and bimodality.}
Replicating bug01 --- required by our pre-registered rule for pairs whose arms
differ by under 30\% --- showed Arm~A is \emph{bimodal}: two runs produced a
minimal one-byte repair (\$8.66, \$11.78, pristine binary) and two produced a
redesign of the quiescence policy (\$44.27, \$41.35). Arm~B produced the minimal
repair in all three runs (\$49.63, \$79.28, \$51.40). Matched on solution, FCDD
costs a median $5.0\times$ (range 4.2--9.2). This is distinct from stochastic
token variance~\cite{bai2026}: it is strategy selection, and the expensive Arm~A
branch is the one graders scored at or below the gate floor --- which the sealed
acceptance script passed regardless.

\paragraph{Quality.}
All submissions cleared the pre-registered threshold. After the first grading
round proved order-confounded, we re-graded all seven pairs in both presentation
orders with 14 independent graders: position came out exactly 7/7, every pair was
order-consistent, and FCDD was favoured on 6 of 7 ($p=0.125$). Axis means
(A / B): correctness 4.36 / 4.50; clarity 4.64 / 4.29; test quality 4.14 / 4.86.
On three defects the two source fixes were byte-identical, so the verdict turned
entirely on verification: FCDD produced better \emph{evidence}, not better fixes.

\section{Threats to Validity}

Seven defects, one codebase, one model. The oracle asymmetry above is the most
consequential flaw. The instrument changed four times; only final-version cells
are compared. Five of seven Arm~B and two of seven Arm~A cells required a resume
after infrastructure interruptions, and summed cost therefore overstates FCDD.
The bounded rule caps review rounds, not invocations, granting Arm~B $1.8\times$
the reviewer invocations. Acceptance scripts were written by the fault-injecting
agent, so a shared misconception is undetectable. Grading is label-blind, not
treatment-blind. The authors evaluate their own method; every withdrawn claim is
published rather than deleted.

\section{Conclusions}

Both methods reliably repaired real assembly defects. FCDD was more expensive on
7 of 7 matched pairs and never amortised; matched on solution the gap widens to
$5.0\times$. We make no claim about cost predictability. Quality was close, with
FCDD ahead on test quality. Most transferably: review-loop termination dominates
measured cost, self-reported costs must be reconciled against transcripts, and
grading order must be counterbalanced --- we failed each of these at least once,
and each failure produced a publishable-looking result that was false.

\begin{credits}
\subsubsection{\ackname} None.
\subsubsection{\discintname}
The method under test was authored by the same human--agent pair conducting this
evaluation. Mitigations: pre-registration, independently seeded faults with a
sealed key, mechanical gates written before repair, blinded counterbalanced
grading, adversarial review against raw data, and publication of every withdrawn
claim. External replication is warranted.
\end{credits}

\bibliographystyle{splncs04}
\bibliography{refs}

\end{document}
